\section{The analytic estimate: one bound for every configuration}
\label{analytic:section}

The real estimate has three steps: express the determinant as a positive
integral, bound its logarithm by an energy comparison, and restore the
factorial normalization. Fauzan supplies the method and comparison
measure \cite[Section 6 and Appendix A]{fauzan}; Firsching's analytic
files give a further implementation \cite{firsching}.
The detailed companion, Sections 5.1--5.6, proves the supporting bounds
and records the finite certificate. Here we keep their role in the
argument visible.
Throughout, \(\alpha=N/K=3/40\) and \(\lambda=h/K=37/40\).

\subsection{From the determinant to logarithmic energy}

Applying Andr\'eief's identity to the positive Gram matrix gives
\begin{equation}\label{analytic:andreief}
 \Delta_K(\xi)=\frac1{h!}\int_{(0,\infty)^h}
 \prod_{i<j}(y_i^2-y_j^2)^2
 \prod_i A_K(y_i^2)w(y_i)\,dy_1\cdots dy_h.
\end{equation}
Expanding the two Vandermonde determinants and integrating term by term
proves the identity: relabeling the variables leaves \(h!\) copies of
the Gram determinant. Exponential decay makes every term integrable
\cite{andreief}.

Put \(y_i=K\sqrt{t_i}\). The external field comes directly from the
finite products:
\[
 \log D_m(K^2t)
 =2m\log K+K\int_0^{m/K}\log(t+u^2)\,du+O(\log K),
\]
uniformly for \(t\ge0\) and \(1\le m\le K\).
The error is between zero and \(2\log K+2\), by comparison with the
right-endpoint Riemann sum (companion, Section 5.5).
Including the exponential factor in \(w\) gives
\begin{equation}\label{analytic:field}
 V(t)=2\pi\sqrt t+\int_0^1\log(t+u^2)\,du
                 -6\int_0^\alpha\log(t+u^2)\,du.
\end{equation}
Thus the main exponent of the integrand is
\[
 E_K(\mathbf t)=2\sum_{i<j}\log|t_i-t_j|-K\sum_i V(t_i).
\]
The first term rewards separation; the second confines the points.
This is the real counterpart of taking valuations of the same products
in the arithmetic estimate.

\subsection{Why one comparison measure is enough}

Write \(U^\rho(t)=\int\log|t-u|\,d\rho(u)\) and
\(I(\rho)=\iint\log|t-u|\,d\rho(t)d\rho(u)\).
Fauzan's explicit positive combination of sixteen arcsine measures
has mass \(\lambda\), support in \((0,2)\), and satisfies
\begin{equation}\label{analytic:potential-target}
 2U^\rho(t)-V(t)\le M_0=-\frac{1329}{200}
 \qquad(t\ge0).
\end{equation}
Its potential is certified on all 684 intervals of a partition of
\([0,2]\), using outward rational bounds; a decreasing elementary
majorant handles \(t\ge2\). The data and verification are in the
companion, Sections 5.2--5.3, and \texttt{checks/analytic\_certificate.py}.
This measure is chosen to prove a useful inequality, not to solve
the exact equilibrium problem.

The reason the inequality works is short. With this \(\rho\), let
\(\eta\) be positive and compactly supported on \([0,\infty)\), of
mass \(\lambda\), with absolute logarithmic integrability for
\((\eta+\rho)\otimes(\eta+\rho)\).
The zero-mass energy inequality gives \(I(\eta-\rho)\le0\), hence
\[
 I(\eta)-\int V\,d\eta
 =I(\eta-\rho)+\int(2U^\rho-V)\,d\eta-I(\rho)
 \le\lambda M_0-I(\rho).
\]
The Gaussian-kernel proof of the zero-mass inequality is in the
companion, Section 5.4. The original moment measure defines the
determinant; \(\eta\) describes a configuration; \(\rho\) bounds all
configurations at once.

The configuration measure \(\eta=K^{-1}\sum_i\delta_{t_i}\) has mass
\(\lambda\). With the diagonal omitted,
\(I(\eta)-\int V\,d\eta=E_K/K^2\).
Its self-energy is infinite and negative, so it cannot be inserted
directly. Replace each point mass by arclength
on a circle of radius \(K^{-2}\). The circle self-energies and the
change in the real-axis potential produce only \(O(K\log K)\).
Keeping an extra decaying factor gives the precise usable bound
\begin{align}\label{analytic:configuration}
 E_K(\mathbf t)+\sum_i\sqrt{t_i}
 \le{}&(\lambda M_0-I(\rho))K^2\\
      &+(120+\sqrt2)h+2h\log K,\notag
\end{align}
for every \(t_i\ge0\) and admissible \(K\).
The companion's Section 5.4 proves the circle-average and tail
estimates; a real-axis potential inequality alone would not suffice.
The retained factor \(\exp(-\sum_i\sqrt{t_i})\) makes the bound integrable.

\subsection{The normalization cancels the larger terms}

The remaining factors are elementary. The bound
\(w(y)\le8192(1+y)^5e^{-2\pi y}\), the Jacobians, and the product
errors contribute \(O(K\log K)\) to the logarithm. After
\eqref{analytic:configuration}, the remaining one-variable integral is
\(\int_0^\infty t^{-1/2}(1+\sqrt t)^5e^{-\sqrt t}\,dt=652\).
The powers of \(K\) in the Vandermonde and \(A_K\) give
\(\kappa=2\lambda(\lambda+6\alpha-1)\), so
\[
 \log\Delta_K(\xi)
 \le\kappa K^2\log K+(\lambda M_0-I(\rho))K^2+O(K\log K).
\]
The factorial bounds for \(S_K\) give the opposite leading term:
\begin{align*}
 \log S_K&\le-\kappa K^2\log K+C_*K^2+O(K\log K),\\
 C_*&=-2\lambda+12\alpha\lambda(1-\log\alpha)
             +3\lambda^2-2\lambda^2\log(2\lambda).
\end{align*}
Thus normalization cancels the \(K^2\log K\) terms exactly. The
certificate verifies
\[
 \lambda M_0-I(\rho)+C_*<-1.3669955.
\]
Keeping the elementary errors explicit yields, for every \(K=40n\),
\begin{equation}\label{analytic:goal}
 \boxed{\ \log F_K(\xi)
 \le-1.3669955K^2+27K\log K+200K.\ }
\end{equation}
Sections 5.5--5.6 of the companion carry out this bookkeeping.
The important margin is the negative quadratic coefficient; the
arithmetic must make the coefficients integral at a smaller
quadratic cost.
